\section{Conclusion}
\label{sec:discussion}

Successfully following a skill does not make its internal states an
independently reusable substitute for the source. Our experiments reveal a
depth interval in which skill-state transfer has weakened while the original
computation continues to depend on skill attention. This separation recurs
across three models and two task domains, including common-item evaluation.
The model can still use a skill after its states have become ineffective at
correcting another execution.

Content and cross-question controls show that effective early transfer
carries structured skill information. Matching state norms preserves that
effect, while repeated late replacement does not restore it. Source-cache repair further shows that late residual-transfer failure is
partly reversible: pairing states with their cached history raises layer-$16$
accuracy from $3.5\%$ to $42.1\%$. Yet cutting subsequent source reading
reduces it to $7.0\%$. Separating attention paths confirms that late
source dependence is driven by continued reading. Useful skill information
can therefore survive a failed residual transplant while requiring a richer
execution state and further access to its source.

Contextual competence therefore has two distinct properties: the ability to
support the current computation and the ability to be reused in another.
For skill-sharing agents and context-reuse systems, both need to be tested.
A representation that works within its original computation is not yet a
portable replacement for the instructions that produced it.
